%! Transformations of Functions — Unit 1, Lesson 1.4 — Group Activity (ANSWER KEY)
\documentclass[10pt]{article}
\usepackage{precalculus-article}
\usepackage{precalculus-key}
\newcommand{\TallMath}[1]{$\displaystyle #1\rule[-1.4em]{0pt}{3.2em}$}

\begin{document}

\pageheader{Unit 1, Lesson 1.4}{Group Activity --- Answer Key}

\vspace{0.12in}

\begin{scenariobox}[The Fountain Deck]{plumlight}
\small
A city fountain is being designed. A measuring tape runs along the deck, and
\textbf{nozzle A} stands at the $0$ mark. The function $f$ gives the height of
nozzle A's water arc, in \textbf{feet}, above the point $x$ feet along the
tape. Every other nozzle on the plan is written as a rule built out of $f$ ---
that is the whole spec sheet.

\begin{center}
\begin{tikzpicture}[x=0.58cm, y=0.16cm]
  \draw[linegray, very thin, xstep=1, ystep=4] (0,0) grid (8,16);
  \draw[->, thick] (0,0) -- (8.8,0) node[below right] {\small $x$ (ft along the deck)};
  \draw[->, thick] (0,0) -- (0,17.8) node[above] {\small $f(x)$ (ft)};
  \foreach \x in {1,2,3,4,5,6,7,8} \node[below, font=\small] at (\x,0) {\x};
  \foreach \y in {4,8,12,16} \node[left, font=\small] at (0,\y) {\y};
  \draw[very thick, plum, smooth, samples=60, domain=0:8]
    plot (\x, {-\x*\x + 8*\x});
  \fill[plum] (0,0) circle (2.5pt);
  \fill[plum] (4,16) circle (2.5pt);
  \fill[plum] (8,0) circle (2.5pt);
\end{tikzpicture}
\end{center}

\renewcommand{\arraystretch}{1.3}
\begin{center}
\begin{tabular}{|c|c|c|c|c|c|c|c|c|c|}
\hline
$x$ (ft) & 0 & 1 & 2 & 3 & 4 & 5 & 6 & 7 & 8 \\
\hline
$f(x)$ (ft) & 0 & 7 & 12 & 15 & 16 & 15 & 12 & 7 & 0 \\
\hline
\end{tabular}
\end{center}
\end{scenariobox}

\vspace{0.1in}

\boxguard[26]
\begin{tcolorbox}[colback=white, colframe=black!40,
    title=\textbf{Tier R --- Remediate},
    fonttitle=\bfseries, arc=2mm, left=3mm, right=3mm, top=2mm, bottom=2mm]

\begin{enumerate}[itemsep=4pt]
  \item Read straight off the table: $f(0) = $ \ans{0} feet (the water
        leaves the nozzle) and $f(4) = $ \ans{16} feet (the top of the arc).

        The high point of nozzle A's arc is at ( \ans{4} ,
        \ans{16} ).

  \item \textbf{Nozzle P} sits on a 3-foot pedestal, so its arc is
        $P(x) = f(x) + 3$.

        \smallskip
        \renewcommand{\arraystretch}{1.4}
        \begin{center}
        \begin{tabular}{|c|c|c|c|c|c|}
        \hline
        $x$ (ft) & 0 & 2 & 4 & 6 & 8 \\
        \hline
        $f(x)$ & 0 & 12 & 16 & 12 & 0 \\
        \hline
        $P(x) = f(x)+3$ & \ans{3} & \ans{15} & \ans{19} & \ans{15} & \ans{3} \\
        \hline
        \end{tabular}
        \end{center}
        \smallskip

        Was the pedestal a change to the \textbf{question} or the
        \ans{answer}? (Circle one.)

  \item \textbf{Nozzle Q} runs on a spec that doubles every height:
        $Q(x) = 2f(x)$.

        \begin{work}
          Q(2) &= 2 \cdot 12 \\
               &= 24
        \end{work}

        \renewcommand{\arraystretch}{1.4}
        \begin{center}
        \begin{tabular}{|c|c|c|c|c|c|}
        \hline
        $x$ (ft) & 0 & 2 & 4 & 6 & 8 \\
        \hline
        $Q(x) = 2f(x)$ & \ans{0} & \ans{24} & \ans{32} & \ans{24} & \ans{0} \\
        \hline
        \end{tabular}
        \end{center}

  \item High points. $P$'s is at ( \ans{4} , \ans{19} ) and $Q$'s is
        at ( \ans{4} , \ans{32} ).

        \vspace{0.05in}
        Both are still at the same $x$ as nozzle A's. Say why in one sentence.

        \vspace{0.05in}
        \ansline{Neither rule changed the question handed to $f$ --- the nozzle never moved along the deck.}
\end{enumerate}
\end{tcolorbox}

\vspace{0.1in}

\boxguard[26]
\begin{tcolorbox}[colback=white, colframe=black!40,
    title=\textbf{Tier A --- Approaching Proficiency},
    fonttitle=\bfseries, arc=2mm, left=3mm, right=3mm, top=2mm, bottom=2mm]

\begin{enumerate}[itemsep=8pt]
  \item \textbf{Nozzle B} is bolted 4 feet farther along the deck. The
        engineer writes its arc as $B(x) = f(x - 4)$.

        \begin{enumerate}[label=(\alph*), itemsep=4pt]
          \item $B(6) = f($ \ans{2} $) = $ \ans{12} feet.
          \item Fill in the rest.

                \smallskip
                \renewcommand{\arraystretch}{1.4}
                \begin{center}
                \begin{tabular}{|c|c|c|c|c|c|}
                \hline
                $x$ (ft) & 4 & 6 & 8 & 10 & 12 \\
                \hline
                handed to $f$: $x - 4$ & \ans{0} & \ans{2} & \ans{4} & \ans{6} & \ans{8} \\
                \hline
                $B(x)$ & \ans{0} & \ans{12} & \ans{16} & \ans{12} & \ans{0} \\
                \hline
                \end{tabular}
                \end{center}
                \smallskip

          \item Compared with nozzle A, $B$'s list of heights is
                \ans{identical}. Only the \ans{positions} changed.
          \item $B$'s high point is at ( \ans{8} , \ans{16} ).
        \end{enumerate}

  \item The nozzle moved to the \textbf{right}, and the rule has a
        \textbf{minus} sign. Explain the mismatch using the words
        \textit{handed to $f$}.

        \vspace{0.05in}
        \ansline{To get nozzle A's top-of-arc height, $f$ must be handed 4. At $x = 8$ the rule hands}
        \ansline{$f$ the value $8 - 4 = 4$, so the peak now happens 4 feet farther right.}

  \item \textbf{Nozzle C} is on the 3-foot pedestal \textit{and} 4 feet down
        the deck: $C(x) = f(x - 4) + 3$. Follow each landmark.

        \smallskip
        \renewcommand{\arraystretch}{1.5}
        \begin{center}
        \begin{tabular}{|c|c|c|c|}
        \hline
        \textbf{point on A's arc} & \textbf{new input} & \textbf{new height} & \textbf{lands at} \\
        \hline
        $(0,\,0)$ & \ans{4} & \ans{3} & \ans{$(4,3)$} \\
        \hline
        $(4,\,16)$ & \ans{8} & \ans{19} & \ans{$(8,19)$} \\
        \hline
        $(8,\,0)$ & \ans{12} & \ans{3} & \ans{$(12,3)$} \\
        \hline
        \end{tabular}
        \end{center}

  \item Nozzle A's water is in the air for $0 \le x \le 8$ and reaches heights
        from $0$ to $16$ feet --- so $f$ has domain $[0, 8]$ and range
        $[0, 16]$.

        \begin{work}
          0 &\le x - 4 \le 8 \\
          4 &\le x \le 12
        \end{work}

        \begin{enumerate}[label=(\alph*), itemsep=4pt]
          \item Domain of $C$: \ans{$[4,\,12]$} \qquad Range of $C$: \ans{$[3,\,19]$}
          \item Which of $C$'s two moves changed where the water lands, and
                which changed how high it goes?

                \vspace{0.05in}
                \ansline{The $x-4$ (a question change) moved the splash zone; the $+3$ (an answer change) raised it.}
        \end{enumerate}
\end{enumerate}
\end{tcolorbox}

\vspace{0.1in}

\boxguard[26]
\begin{tcolorbox}[colback=white, colframe=black!40,
    title=\textbf{Tier E --- Extension},
    fonttitle=\bfseries, arc=2mm, left=3mm, right=3mm, top=2mm, bottom=2mm]

\begin{enumerate}[itemsep=8pt]
  \item \textbf{Reverse-engineer the spec.} A mystery nozzle $M$ was measured
        on site.

        \smallskip
        \renewcommand{\arraystretch}{1.4}
        \begin{center}
        \begin{tabular}{|c|c|c|c|c|c|}
        \hline
        $x$ (ft) & $-3$ & $-1$ & 1 & 3 & 5 \\
        \hline
        $M(x)$ (ft) & 0 & 12 & 16 & 12 & 0 \\
        \hline
        \end{tabular}
        \end{center}
        \smallskip

        \begin{enumerate}[label=(\alph*), itemsep=4pt]
          \item $M(x) = f(x + $ \ans{3} $)$
          \item In plain words, the nozzle was moved \ans{3} feet
                \ans{backward} along the deck from A.
        \end{enumerate}

  \item \textbf{The reflecting pool.} The deck surface is still water, and it
        mirrors the arc below the deck: $R(x) = -f(x)$.

        \smallskip
        \renewcommand{\arraystretch}{1.4}
        \begin{center}
        \begin{tabular}{|c|c|c|c|c|c|}
        \hline
        $x$ (ft) & 0 & 2 & 4 & 6 & 8 \\
        \hline
        $R(x) = -f(x)$ & \ans{0} & \ans{$-12$} & \ans{$-16$} & \ans{$-12$} & \ans{0} \\
        \hline
        \end{tabular}
        \end{center}
        \smallskip

        \begin{enumerate}[label=(\alph*), itemsep=4pt]
          \item Question change or answer change? \ans{answer}
          \item A second reflecting pool 4 feet down the deck shows
                $-f(x - 4)$. Its lowest point is at ( \ans{8} ,
                \ans{$-16$} ).
        \end{enumerate}

  \item \textbf{The scale model versus the double-pressure spec.} A tabletop
        model reproduces the arc in half the deck distance: $m(x) = f(2x)$.
        Compare it against $Q(x) = 2f(x)$ from Tier R.

        \smallskip
        \renewcommand{\arraystretch}{1.4}
        \begin{center}
        \begin{tabular}{|c|c|c|c|c|c|}
        \hline
        $x$ (ft) & 0 & 1 & 2 & 3 & 4 \\
        \hline
        $m(x) = f(2x)$ & \ans{0} & \ans{12} & \ans{16} & \ans{12} & \ans{0} \\
        \hline
        $Q(x) = 2f(x)$ & \ans{0} & \ans{14} & \ans{24} & \ans{30} & \ans{32} \\
        \hline
        \end{tabular}
        \end{center}
        \smallskip

        \begin{enumerate}[label=(\alph*), itemsep=4pt]
          \item The two rules agree at exactly one input: $x = $
                \ans{0}.
          \item $m$'s high point is ( \ans{2} , \ans{16} );
                $Q$'s is ( \ans{4} , \ans{32} ).
          \item The $2$ sits in a different place in each rule. Say what each
                one changed.

                \vspace{0.05in}
                \ansline{$m$'s 2 is inside: it changed the question, squeezing the arc into half the deck.}
                \ansline{$Q$'s 2 is outside: it changed the answer, doubling every height in place.}
        \end{enumerate}
\end{enumerate}
\end{tcolorbox}

\end{document}
